\section{Model YAML Configurations by State Type}
\label{sec:configs}

We present model YAML configuration files for one representative state from each type. These configurations are designed to satisfy the legal criteria of the representative state while remaining within the algorithmic parameter space. Each configuration is associated with the Rust \texttt{redist state} command, which accepts a \texttt{--config} flag pointing to a YAML file.

\subsection{Type I Configuration: Texas (Permissive)}

Texas imposes minimal criteria beyond the federal floor: population equality, contiguity, and VRA compliance for majority-minority opportunities. No compactness requirement is constitutionally mandated, and the legislature retains full discretion over the process.

\begin{verbatim}
# texas_congressional.yaml
# Type I: Permissive
state: TX
year: 2020
districts: 38
resolution: tract
algorithm:
  partition_mode: recursive_bisection
  population_tolerance: 0.005
  compactness: 0.3          # Low: not legally required
  county_weight: 1.0        # Default: no preservation mandate
  coi_weights: null         # No community-of-interest criterion
  partisan_neutral: true    # Structural: no partisan data available
  vra_mode: auto            # Auto-detect majority-minority opportunities
constraints:
  contiguity: required      # Federal floor
  population_balance: strict
output:
  metrics: [polsby_popper, population_deviation, county_splits]
\end{verbatim}

\subsection{Type II Configuration: Iowa (Criteria-Based)}

Iowa is the paradigm case for criteria-based redistricting. The Legislative Services Bureau administers specific statutory criteria: no partisan data may be used, districts must be compact, and county splits must be minimized. Iowa's criteria closely mirror the algorithmic approach, making it the state most naturally suited to algorithmic adoption.

\begin{verbatim}
# iowa_congressional.yaml
# Type II: Criteria-Based (Iowa LSB criteria)
state: IA
year: 2020
districts: 4
resolution: tract
algorithm:
  partition_mode: recursive_bisection
  population_tolerance: 0.005
  compactness: 0.85         # Iowa statute: maximize compactness
  county_weight: 2.0        # Iowa statute: minimize county splits
  coi_weights: null         # Formal communities not specified
  partisan_neutral: true    # Iowa statute: no partisan data
  vra_mode: auto
constraints:
  contiguity: required
  population_balance: strict
  county_splits: minimize   # Primary statutory criterion
output:
  metrics: [polsby_popper, county_splits, population_deviation]
  report_county_splits: true
\end{verbatim}

\subsection{Type III Configuration: California (Commission)}

California's Citizens Redistricting Commission operates under Article XXI of the California Constitution, which mandates criteria in priority order: (1) population equality, (2) VRA compliance, (3) geographical contiguity, (4) community of interest preservation, (5) geographic compactness, (6) city and county boundary respect, (7) nesting of state senate and assembly districts. Proportionality is required as a long-run objective.

\begin{verbatim}
# california_congressional.yaml
# Type III: Commission (California CRC criteria)
state: CA
year: 2020
districts: 52
resolution: tract
algorithm:
  partition_mode: recursive_bisection
  population_tolerance: 0.005
  compactness: 0.85         # CRC criterion #5: geographic compactness
  county_weight: 1.5        # CRC criterion #6: city/county boundaries
  coi_weights: census_places # CRC criterion #4: communities of interest
  partisan_neutral: true    # Structural; proportionality assessed ex post
  vra_mode: strict          # CRC criterion #2: VRA compliance priority
constraints:
  contiguity: required
  population_balance: strict
  vra_compliance: required  # Highest-priority adjustable constraint
output:
  metrics: [polsby_popper, county_splits, majority_minority_districts,
            population_deviation]
  report_vra: true
  report_communities: true
\end{verbatim}

\subsection{Type IV Configuration: Pennsylvania (Hybrid)}

Pennsylvania's redistricting standards were established by the state supreme court in \textit{League of Women Voters v.\ Commonwealth}, which imposed criteria derived from the state constitution's free and equal elections clause: contiguity, compactness (the court's own map achieved mean Polsby-Popper of 0.47), population equality, and avoidance of county and municipality splits where possible.

\begin{verbatim}
# pennsylvania_congressional.yaml
# Type IV: Hybrid (LWV court-imposed standards)
state: PA
year: 2020
districts: 17
resolution: tract
algorithm:
  partition_mode: recursive_bisection
  population_tolerance: 0.001    # LWV standard: near-perfect equality
  compactness: 0.90              # LWV court map: PP mean ~0.47
  county_weight: 1.8             # LWV criterion: minimize county splits
  coi_weights: census_places     # Municipality integrity
  partisan_neutral: true         # Structural
  vra_mode: auto
constraints:
  contiguity: required
  population_balance: strict
  minimum_pp: 0.35               # LWV benchmark threshold
output:
  metrics: [polsby_popper, county_splits, municipality_splits,
            population_deviation]
  report_lwv_compliance: true    # Flag districts below LWV benchmark
\end{verbatim}

\subsection{Type V Configuration: Florida (Restrictive)}

Florida's Amendment 6 imposes the strictest criteria of any state in this analysis. Districts must not be drawn with intent to favor or disfavor a political party or incumbent. Additionally, Article III, Section 20 requires that districts, where feasible, utilize existing political and geographical boundaries, and be as equal in population as practicable, be compact, and where feasible be contiguous. The algorithmic approach satisfies the partisan-neutrality requirement structurally and the compactness and boundary requirements through parameter configuration.

\begin{verbatim}
# florida_congressional.yaml
# Type V: Restrictive (Amendment 6 / Article III Sec. 20)
state: FL
year: 2020
districts: 28
resolution: tract
algorithm:
  partition_mode: recursive_bisection
  population_tolerance: 0.005
  compactness: 0.85         # Art. III Sec. 20: compact districts
  county_weight: 1.5        # Art. III Sec. 20: existing boundaries
  coi_weights: census_places
  partisan_neutral: true    # Amend. 6: structural satisfaction
  incumbent_blind: true     # Amend. 6: no incumbent address data
  vra_mode: strict          # Section 2 compliance; Art. III Sec. 20
constraints:
  contiguity: required
  population_balance: strict
  partisan_data: forbidden  # Constitutional prohibition
  incumbent_data: forbidden # Constitutional prohibition
output:
  metrics: [polsby_popper, county_splits, majority_minority_districts,
            population_deviation]
  certify_partisan_blind: true    # Audit trail for constitutional compliance
  certify_incumbent_blind: true
\end{verbatim}

The Florida configuration is notable for the \texttt{certify\_partisan\_blind} and \texttt{certify\_incumbent\_blind} flags, which generate audit trail records documenting that no partisan or incumbent data was provided to the algorithm at any point during the run. This documentation is valuable for demonstrating constitutional compliance under Amendment 6 in litigation.
